\documentclass[11pt]{article}
\usepackage[a4paper,margin=1.9cm]{geometry}
\usepackage[T1]{fontenc}
\usepackage{textcomp}
\usepackage[spanish]{babel}
\usepackage{amsmath,amssymb}
\usepackage{lmodern}
\renewcommand{\familydefault}{\sfdefault}
\usepackage{enumitem}

\newcommand{\vect}[2]{\begin{pmatrix}#1\\#2\end{pmatrix}}
\usepackage{tikz}
\usetikzlibrary{calc,arrows.meta,patterns,decorations.pathmorphing}
\usepackage[table]{xcolor}
\usepackage{multicol}
\usepackage{parskip}

\definecolor{unalblue}{RGB}{31,73,124}

\setlist[enumerate,1]{itemsep=1.4em,topsep=1em}

\newcommand{\examheader}{%
  \noindent
  \begin{minipage}[t]{0.6\linewidth}
    {\small\bfseries Geometría Vectorial y Analítica --- Parcial 1}\\
    {\small Semestre 2014--02}
  \end{minipage}%
  \hfill
  \begin{minipage}[t]{0.35\linewidth}
    \raggedleft\small Escuela de Matemáticas. Universidad Nacional de Colombia.
  \end{minipage}\\[2pt]
  \rule{\linewidth}{0.6pt}
}
\newcommand{\examfooter}{%
  \vfill
  \rule{\linewidth}{0.4pt}\\[1pt]
  {\scriptsize\textit{Transcripción digital --- MathUNAL}\hfill\textcolor{unalblue}{\textbf{\thepage}}}
}
\pagestyle{empty}

\newcommand{\nota}[1]{{\small\itshape\color{unalblue!70!black} #1}}

\begin{document}
\examheader

\begin{center}
{\Large Primer Examen Parcial de Geometría Vectorial}\\[4pt]
{\small 8 de septiembre de 2014}
\end{center}

\vspace{10pt}
\textbf{PUNTAJE. SOLO PARA USO OFICIAL}

\begin{center}
\begin{tabular}{|c|c|c|c|c|c|}
\hline
1 & 2 & 3 & 4 & 5 & TOTAL \\
\hline
 & & & & & \\
\hline
\end{tabular}
\end{center}

\vspace{10pt}
\textbf{INSTRUCCIONES}

No está permitido el uso de calculadora ni de otros aparatos electrónicos. Solucione las preguntas en los espacios en blanco asignados a cada una de ellas. \textbf{Duración: 1 hora y 50 minutos.} No está permitido sacar hojas en blanco ni ningún tipo de apuntes durante el examen. Utilice la página 6 para realizar sus operaciones en borrador.

\begin{enumerate}
\item (12\%) Considere las rectas $\mathcal{L}_1: (2s)x+y-3=0$ y $\mathcal{L}_2: (5s)x+y+2=0$ con $s\in\mathbb{R}$. ¿Hay algún valor de $s$ para el cual el ángulo dirigido de la recta $\mathcal{L}_1$ a la recta $\mathcal{L}_2$ es de $45^\circ$? Justifique su respuesta.

\item Sean $A=\vect{2}{-3}$ y $B=\vect{-1}{5}$.

\begin{enumerate}[label=(\alph*)]
\item (12\%) Muestre que $A$ y $B$ son linealmente independientes.
\item (14\%) Encuentre vectores $P$ y $Q$ tales que $B=P+Q$ siendo $P$ múltiplo escalar de $A$ y $Q$ ortogonal a $A$.
\end{enumerate}

\item \

\begin{enumerate}[label=(\alph*)]
\item (14\%) Muestre que la recta $\mathcal{L}$ con ecuaciones escalares paramétricas $x=2-5t$, $y=5+2t$ es paralela a la recta $\mathcal{L}'$ con ecuación general $4x+10y-13=0$.
\item (12\%) Halle la distancia entre las rectas $\mathcal{L}$ y $\mathcal{L}'$ del literal (a).
\end{enumerate}

\item Considere la recta $\mathcal{L}$ con ecuación general $\sqrt{3}x+y-7=0$ y la recta $\mathcal{L}'$ que pasa por el punto $\vect{-1}{1}$, tal que el ángulo dirigido de $\mathcal{L}$ a $\mathcal{L}'$ es $30^\circ$.

\begin{enumerate}[label=(\alph*)]
\item (12\%) Halle una ecuación en la forma vectorial para la recta $\mathcal{L}'$.
\item (12\%) Encuentre una ecuación en la forma general para la recta que pasa por el punto $\vect{-1}{1}$ y es perpendicular a la recta $\mathcal{L}$.
\end{enumerate}

\item (12\%) Sea $\mathcal{C}$ un cuadrilátero con vértices $A$, $B$, $C$ y $D$. Pruebe que los puntos medios de los lados del cuadrilátero $\mathcal{C}$ son los vértices de un paralelogramo.

\end{enumerate}

\examfooter
\end{document}
